Las diferentes maneras de representaciones digitales que existen para el sonido digital no solo son una secuencia de muestras, sino que además se guardan en diferentes formatos:

\begin{itemize}
    \item \textbf{Formatos sin compresión:} \textit{WAV, AIFF}. Formatos que guardan directamente las muestras de audio codificadas, ofrecen mayor fidelidad al original pero también generan archivos más pesados. Es el formato utilizado en todo el proyecto.
    \item \textbf{Formatos comprimidos sin pérdida:} \textit{FLAC, ALAC}. Formatos que reducen el tamaño sin alterar la información.
    \item \textbf{Formatos comprimidos con pérdida:} \textit{MP3, AAC, OGG...}. Utilizan modelos psicoacústicos para descartar información que el oído humano no percibe. Se usan en transmisión y almacenamiento masivo de audio para consumo (\textit{streaming}, etc.).
\end{itemize}
